\subsection{Whole-archive refitting, retention and uncertainty}
The whole-archive benchmark comprises \WARepetitions{} station/repetition/setting evaluations at \WAStations{} long-record stations in 18 connected groups. Figure~\ref{fig:whole-archive} and Table~\ref{tab:whole-archive-paired} use common station/repetition rows. Retaining years improves the finite-record slope comparison in the evaluated settings. The ranking of D2 and calibrated stochastic regression varies with time setting and region; results are conditional on the donor availability and mask design. Twenty masks at one station are not twenty independent rivers or climates.

The factorial comparison identifies a retention contribution rather than assigning one step of a sequential pipeline all of its improvement. For offline controlled gaps, the symmetric retention contribution is \WARetentionShare\% of the D2-versus-rule/plug-in MSE gain. This is a result for this experiment, not a universal percentage. Estimator effects and interactions are retained in the full paired table.

Annual marginals and trend uncertainty give different answers (Figure~\ref{fig:joint-uncertainty}). Under controlled gaps the D2 missing-data trend interval covers \WAOfflineCoverage\% offline and \WAForwardCoverage\% forward; residual blocks give \WABlockForwardCoverage\% forward. The tail and dependence sensitivities are separated: the same bounded annual marginals are used for each rank coupling, whereas endpoint and normal-tail assumptions are compared on paired cases. Annual miscalibration, extrapolated tails and dependence can all contribute to trend errors; the experiment does not uniquely attribute the residual undercoverage to any one of them. Task calibration is available for \WATaskForwardN{} of 820 forward controlled-gap D2 cases (\WATaskForwardPercent\%), representing two stations; both task calibration and the independent-draw comparator cover all of these same cases. This selected-subset result does not demonstrate general improvement. Withholding and regional support remain explicit.

\subsection{Observed facts, numerical estimates and unknown outputs}
The release logic is stated in Figure~\ref{fig:release-system}. It is a design schematic; the empirical support for each gate is reported below and does not constitute an operational certification.
The three-output comparison uses a point estimate only within the source feature range and a raw individual classifier-score cap, an observed-window bound when a numerical estimate is withheld, and unknown when no valid bound is available or visible readings raise a quality flag. The classifier was trained for the original unprojected D2 and is not a calibrated group-risk bound for the projected estimate. Visible anomaly limits come from the corresponding source training split, including pre-2010 data for forward evaluation. A seven-day water--air agreement flag addresses the specified exact-air-substitution scenario only; it is not evaluated in the other stress scenarios and is not a general sensor diagnostic. Where a valid seven-day water--air window is unavailable, this diagnostic is unassessed and does not itself block release. The primary factorial trend comparison is separate; an additional exploratory analysis passes released values and bounds to a Gaussian censored working likelihood.

Figure~\ref{fig:output-boundaries} reports the cost of withholding numbers. Source-range diagnostics do not make hot-year extrapolation or temporal shift safe. Exact-air substitution can be detected while sacrificing almost all numerical output in that condition; zero publication has undefined selected risk and is not successful numerical recovery. Observed bounds retain heat information only when the reported measurements are valid. Event tables distinguish detected hot-year cases, missed or withheld cases, and false hot declarations; mask realizations are not counted as independent physical events. Mixture curves vary the imposed contamination fraction; they do not estimate its prevalence in real archives.

In the exploratory release-then-trend analysis, the forward controlled-gap RMSE is \WASelectedRMSE{}$^{\circ}$C per decade when selected values and censored bounds enter the Gaussian working likelihood, compared with \WASelectedBaseRMSE{} for retaining all D2 estimates on the same cases. The nominal 95\% working interval covers \WASelectedCoverage\%, but its average width is \WASelectedWidth{}$^{\circ}$C per decade. That interval includes fitted regression variation, treats selected numerical estimates as exact and is not a calibrated missingness-only interval. Selection, informative censoring and the working distribution can all affect the result. In offline controlled gaps, \WAQualityUnassessed{} bound cases have no evaluable air-agreement diagnostic. Withholding those bounds leaves numerical publication unchanged and gives a working-likelihood RMSE of \WAQualityStrictRMSE{}; this sensitivity does not establish sensor safety.
